\subsection{Generic groups}

So far we have treated the Fourier transform on the cyclic group $\cyclicgroupof{\prod_{\catenumeratorin}\catdimof{\catenumerator}}$ using unfoldings.
More naturally, a tensor is defined as a complex field on the cartesian product group $\facstates$.
Let us now generalize to generic groups.

\subsubsection{Representations}

Recall that tensors are group fields, see \secref{sec:groupFields}.

A representation is a map of $\gengroup$ to a matrix group.
Irreducible representations (irrep) are those, which are not equivalent to matrix groups in block form.
We enumerate them by a variable $\headvariable$ and store all irreps in the tensor
\begin{align*}
    \repofat{\gengroup}{
        \catvariable,\headvariable,\cselvariable{\insymbol},\cselvariable{\outsymbol}
    } \, ,
\end{align*}
where the $\headindex$-th representation maps the group element $\catindex$-th to the matrix
\begin{align*}
    \frac{1}{\sqrt{\seldim}} \cdot \repofat{\gengroup}{
        \indexedcatvariable,\indexedheadvariable,\cselvariable{\insymbol},\cselvariable{\outsymbol}
    } \, .
\end{align*}
Here $\seldim$ is the dimension of the representation.
We here made the convention to store in the slices of $\repof{\gengroup}$ the by ${\sqrt{\seldim}}$ rescaled irrep, to ease the definition of the Fourier transform.




\begin{example}{Cyclic group}\label{exa:cycGroupIrreps}
    For the cyclic group $\gengroup=\cyclicgroupof{\catdim}$ the irreps are one-dimensional, that is $\headdim=1$, and enumerated by $\headindex\in[\catdim]$ as
    \begin{align*}
        \repofat{\cyclicgroupof{\catdim}}{\indexedcatvariable,\indexedheadvariable,\cselvariable{\insymbol}=0,\cselvariable{\outsymbol}=0}
        = \expof{\twopii \cdot \frac{\catindex \cdot \headindex}{\catdim}} \, .
    \end{align*}
    The characters of the irreps are
    \begin{align*}
        \characterofat{\headindex}{\catindex}=\repofat{\cyclicgroupof{\catdim}}{\indexedcatvariable,\indexedheadvariable,\cselvariable{\insymbol}=0,\cselvariable{\outsymbol}=0} \, .
    \end{align*}
    Note that for $\catdim=2$ this is the Hadamard matrix
    \begin{align*}
        \repofat{\cyclicgroupof{2}}{\catvariable,\headvariable,\cselvariable{\insymbol}=0,\cselvariable{\outsymbol}=0}
        = \sqrt{2} \cdot \hgateat{\catvariable,\headvariable} \, .
    \end{align*}
\end{example}

\subsubsection{Finite Abelian groups}

Finite Abelian groups are equivalent to cartesian products of cyclic groups, that is
\begin{align*}
    \gengroup \sim \facgroup \, .
\end{align*}
The irreps of finite Abelian groups are thus disconnected tensor networks
\begin{align*}
    \repofat{\facgroup}{\shortcatvariables,\headvariableof{[\catorder]},\cselvariableof{\insymbol}{[\catorder]},\cselvariableof{\outsymbol}{[\catorder]}}
    = \bigotimes_{\catenumeratorin} \repofat{\cyclicgroupof{\catdimof{\catenumerator}}}{
        \catvariableof{\catenumerator},\headvariableof{\catenumerator},\cselvariableof{\insymbol}{\catenumerator},\cselvariableof{\outsymbol}{\catenumerator}
    }
\end{align*}
where each $\repofat{\cyclicgroupof{\catdimof{\catenumerator}}}{\catvariableof{\catenumerator},\headvariableof{\catenumerator},\cselvariableof{\insymbol}{\catenumerator},\cselvariableof{\outsymbol}{\catenumerator}}$ is the irrep tensor to $\cyclicgroupof{\catdimof{\catenumerator}}$, see \exaref{exa:cycGroupIrreps}.

\subsubsection{Fourier transform}

The Fourier transform of a complex group field $\hypercoreat{\catvariable}$ is the tensor
\begin{align*}
    \sqrt{\frac{1}{\cardof{\gengroup}}} \cdot \contractionof{
        \hypercoreat{\catvariable},\repofat{\gengroup}{\catvariable,\headvariable,\cselvariable{\insymbol},\cselvariable{\outsymbol}}
    }{\headvariable,\cselvariable{\insymbol},\cselvariable{\outsymbol}} \, .
\end{align*}

\begin{example}[Walsh-Hadamard transform]
    Consider the group $\gengroup=\atomgroup$, and a complex field on it by a quantum state $\qstateat{\shortcatvariables}$.
    The Fourier transform is then the application of Hadamard gates on each qubit, that is
    \begin{align*}
        \contractionof{
            \{\qstateat{\shortcatvariables}\}
            \cup \bigcup_{\catenumeratorin}
            \{\hgateat{\catvariableof{\catenumerator},\headvariableof{\catenumerator}}\}
        }{\headvariableof{[\catorder]}} \, .
    \end{align*}
    This is the Walsh-Hadamard transform of the quantum state.
\end{example}

\subsubsection{Coset transform}

Let us again assume finite Abelian groups $\gengroup=\facgroup$.
Applying the Fourier transform on a coset encoding tensor gives a tensor supported at the annihilator.

\begin{lemma}\label{lem:cosetTransform}
    Let $\gengroup=\facgroup$, $\gensubgroup\subset\gengroup$ be a subgroup and $\thirdcatindexof{[\catorder]}\gensubgroup$ be the coset to an element $\thirdcatindexof{[\catorder]}$.
    Then
    \begin{align*}
        \contractionof{
            \ftat{\headvariableof{[\catorder]},\shortcatvariables},
            \onehotmapofat{\gencoset}{\shortcatvariables}
        }{\headvariableof{[\catorder]}}
    \end{align*}
    is supported at the elements of $\genannihilator$.
\end{lemma}
\begin{proof}
    We have
    \begin{align*}
        \contractionof{\ftat{\headvariableof{[\catorder]},\shortcatvariables},\onehotmapofat{\gencoset}{\shortcatvariables}}{\indexedheadvariableof{[\catorder]}}
        & = \frac{1}{\sqrt{\cardof{\gengroup}}} \sum_{\shortcatindices\in\gensubgroup} \chi_{\headindexof{[\catorder]}}(\thirdcatindexof{[\catorder]}\otimes \shortcatindices) \\
        & = \frac{1}{\sqrt{\cardof{\gengroup}}} \characterofat{\headindexof{[\catorder]}}{\thirdcatindexof{[\catorder]}} \cdot \left(\sum_{\shortcatindices\in\gensubgroup} \characterofat{\headindexof{[\catorder]}}{\shortcatindices}\right) \, .
    \end{align*}
    Now, if $\chi_{\headindexof{[\catorder]}}\in \genannihilator$, each summand is $1$, and the coordinate evaluates to $\frac{\cardof{\gensubgroup}}{\sqrt{\cardof{\gengroup}}} \cdot \chi_{\headindexof{[\catorder]}}(\thirdcatindexof{[\catorder]})$.
    Using that $\cardof{\gensubgroup}\cdot\cardof{\genannihilator}=\cardof{\gengroup}$, the absolute of the coordinate is
    \begin{align*}
        \sqrt{\cardof{\gensubgroup}} \cdot \frac{1}{\sqrt{\cardof{\genannihilator}}} \, .
    \end{align*}
    The Hilbert-Schmidt norm of the contraction restricted to $\genannihilator$ is therefore $\sqrt{\cardof{\gensubgroup}}$ and coincides with the Hilbert-Schmidt norm of the coset encoding.
    All coordinates to $\hat{\gengroup}\backslash \genannihilator$ have thus to vanish, which shows the claim.
    % Alternative derivation
    Alternatively one can derive the claim, based on destructive inference on the sum term.
\end{proof}

\begin{example}\label{exa:cosetWHTransform}
    Let us consider the group $\gengroup=\atomstates$, which coset representing hypercores are characterized in \exaref{exa:cosetsCyclicTwoProduct}.
    If $\cardof{\genannihilator}=2$ we represent a coset by
    \begin{align*}
        \bigoplus[\catvariableof{\edge}] \quad\text{or}\quad
        \lnot\bigoplus[\catvariableof{\edge}]
    \end{align*}
    where $\edge$ is the support of the generator of $\genannihilator$.
    Notice that
    \begin{align*}
        \bigoplus[\catvariableof{\edge}]
        &= \sqrt{2^{\catorder}} \cdot \left(\bigotimes_{\catenumerator\in\edge} \hplusbasat{\catvariableof{\catenumerator}} + \bigotimes_{\catenumerator\in\edge} \hminusbasat{\catvariableof{\catenumerator}}\right) \\
        \lnot\bigoplus[\catvariableof{\edge}]
        &= \sqrt{2^{\catorder}} \cdot \left(\bigotimes_{\catenumerator\in\edge} \hplusbasat{\catvariableof{\catenumerator}} - \bigotimes_{\catenumerator\in\edge} \hminusbasat{\catvariableof{\catenumerator}}\right)
    \end{align*}
    The Fourier transform for this group is the Walsh-Hadamard transform which turns the basis legs $\hplusbasat{\catvariableof{\catenumerator}}$ and $\hminusbasat{\catvariableof{\catenumerator}}$ to $\fbasisat{\catvariableof{\catenumerator}}$ respectively $\tbasisat{\catvariableof{\catenumerator}}$.
    The Walsh-Hadamard transform therefore returns
    \begin{align*}
        \sqrt{2^{\catorder}} \left(\onehotmapofat{\zerotuple{\edge}}{\catvariableof{\edge}} \pm \onehotmapofat{\onetuple{\edge}}{\catvariableof{\edge}} \right) \, ,
    \end{align*}
    which is supported at $\genannihilator$.
\end{example}

\subsubsection{Hidden subgroup problem}\label{sec:ccHSP}

Lets assume that a function $\exformula:\gengroup\rightarrow \mathcal{Y}$ mapping from a finite abelian group $\gengroup=\facgroup$ to an arbitrary set $\mathcal{Y}$ is constant on the cosets of a subgroup $\gensubgroup$ (called the hidden subgroup), and different on each coset.
The slices
\begin{align*}
    \bencodingofat{\exformula}{\indexedheadvariable,\shortcatvariables}
\end{align*}
of the basis encoding to any image element enumerated by $\headindex$ are then coset encoding tensors and by \lemref{lem:cosetTransform}
\begin{align*}
    \contractionof{
        \bencodingofat{\exformula}{\indexedheadvariable,\shortcatvariables},
        \ftat{\headvariableof{[\catorder]},\shortcatvariables}
    }{\headvariableof{[\catorder]}}
\end{align*}
is supported on the annihilator $\genannihilator$, where each coordinate has equal support.

When retrieving the support of this transform, we can construct a generating set for $\genannihilator$ and then $\gensubgroup$ as $\big(\genannihilator\big)^{\annihilatorsymbol}$ based on the solution of a equation system in modulus calculus.
This strategy is exploited in quantum computing and we will discuss it in \secref{sec:qcHSP}.